\documentclass[11pt,a4paper]{article}
\usepackage[margin=2.4cm]{geometry}
\usepackage{fontspec}
\setmainfont{Times New Roman}
\setmonofont{Menlo}[Scale=0.88]
\usepackage{amsmath,amssymb}
\usepackage{graphicx}
\usepackage{booktabs}
\usepackage{array}
\usepackage[numbers,sort&compress]{natbib}
\usepackage[font=small,labelfont=bf]{caption}
\usepackage{float}
\usepackage[hidelinks]{hyperref}
\usepackage{xcolor}
\newcommand{\todo}[1]{\textcolor{red}{[TODO: #1]}}
\newcommand{\best}[1]{\textbf{#1}}
\newcommand{\worst}[1]{\textbf{#1}}

\title{\bfseries Clean Accuracy Is Not Diagnosis:\\
A Noise-Stress Re-evaluation of Lightweight Bearing Fault Diagnosis}
\author{Anonymous Author(s)\thanks{Manuscript prepared for submission; author names, affiliation and repository URL to be completed at submission. Code, split files and per-seed results accompany the manuscript.}}
\date{}

\setlength{\emergencystretch}{2.5em}
\tolerance=1500
\setcounter{topnumber}{3}
\setcounter{bottomnumber}{2}
\setcounter{totalnumber}{5}
\renewcommand{\topfraction}{0.92}
\renewcommand{\bottomfraction}{0.85}
\renewcommand{\textfraction}{0.06}
\renewcommand{\floatpagefraction}{0.68}

\begin{document}
\maketitle

\begin{abstract}
\noindent
Lightweight one-dimensional convolutional networks are the default choice for bearing fault diagnosis on resource-constrained hardware, and the public CWRU benchmark is their standard proving ground. Reported accuracies have approached 100\% for years, which makes the benchmark increasingly uninformative: if every architecture is saturated, it can no longer separate a robust model from a fragile one. We re-evaluate eight representative lightweight architectures (WDCNN, TICNN, MCNN, ResNet1D, DS-CNN, LeNet1D, ViT1D and MLP-Mixer1D) under a leakage-safe protocol in which no sampling point is shared between training and test windows. Five observations emerge. First, clean-set accuracy is statistically indistinguishable across architectures (97--100\%) and the accuracy inflation caused by overlapping-window leakage is small ($\leq 0.6$ points at full data). Second, the same models differ by up to \best{55 points} when the test signal is corrupted by additive noise at 10\,dB SNR (0.43--0.98), and the clean ranking does not predict the noisy ranking. Third, most of this brittleness is a property of the training recipe rather than of the architecture: per-batch noise augmentation lifts the mean 10\,dB accuracy from 0.70 to 0.97 (Wilcoxon signed-rank, $p<10^{-6}$) and the weakest model from 0.43 to 0.95, at a cost of 5.4 points on clean data. Fourth, a multi-channel wide first-layer stem recovers +13 to +52 points at 10\,dB for about 1\% additional parameters and less than 3 points of clean accuracy, and the number of front-end channels rather than the kernel width is the decisive design knob. Finally, open-set rejection is architecture-dependent and \emph{anti}-correlated with noise robustness (Spearman $\rho=-0.64$): the model that survives noise best rejects unknown faults worst (AUROC 0.84, TPR at 5\% FPR 0.12), a failure that closed-set accuracy tables cannot reveal. The augmentation that repairs noise accuracy does not repair rejection---it lowers mean AUROC from 0.90 to 0.80---so the two axes trade off rather than overlap. The collapse and its repair reproduce on a second rig (Paderborn), and the ordering of the models changes again under harmonic, impulsive and drift noise, so a robustness claim that does not name the perturbation is unfalsifiable. We release the protocol, the split files and every per-seed result.

\vspace{0.5em}
\noindent\textbf{Keywords:} bearing fault diagnosis; lightweight neural networks; noise robustness; leakage-safe evaluation; open-set recognition; reproducibility
\end{abstract}

\section{Introduction}

Rolling-element bearing faults are among the most common causes of failure in rotating machinery, and vibration-based diagnosis is the standard non-invasive monitoring modality \citep{randall2011tutorial}. Over the past decade the field has converged on deep learning: a one-dimensional convolutional network (1D-CNN) takes a short window of the accelerometer signal and outputs a fault class. Because the models and the datasets are both small, the literature has developed a distinctive style: a new block, attention module or multi-scale fusion path is proposed, trained on CWRU, and reported with a table of accuracies that now routinely reach 99\% or even 100\% \citep{zhang2017wdcnn,zhao2019mcnn,m2024a}.

This style of reporting has a structural problem that is easy to overlook from the inside: \emph{a benchmark on which every method scores 99\% cannot be used to choose a method}. If WDCNN reaches 99.7\% and a five-layer ResNet reaches 99.96\%, the difference is smaller than the seed-to-seed variation, yet it is routinely reported as an improvement. Worse, the numbers are not comparable across papers, because split protocols differ: some authors split overlapping windows at random, some split by recording, and some do not say \citep{phme2026,ssrn2024}. Recent position papers have flagged the resulting evaluation problems, but to our knowledge the size of the effect has not been quantified across a common set of architectures under a single harness.

The same literature reports a second, largely separate claim: that the proposed models are “robust to noise”. Robustness is typically demonstrated by adding white Gaussian noise at a handful of signal-to-noise ratios (SNRs) and showing that the new model degrades more gracefully than its predecessors \citep{zhang2018ticnn,mst2024b,sr2025a}. Because the SNRs, the noise model and the training procedure vary from paper to paper, these claims are equally hard to compare.

Our starting point was the suspicion that these two problems are the same problem. To test it we built one evaluation harness in which eight lightweight architectures---chosen to span the design space of the field, from a 3k-parameter LeNet to a 750k-parameter MLP-Mixer---are trained with identical hyperparameters and evaluated on identical splits under (i) four split protocols, (ii) six test SNRs, (iii) with and without noise augmentation during training, and (iv) an open-set setting in which three fault classes are withheld from training. A second dataset (Paderborn) checks that the conclusions are not an artefact of CWRU.

The results were not what we expected. The leakage concern turned out to be quantitatively minor: leaky and leakage-safe splits differ by 0.3 points when the full training set is used. What is not minor is noise: at 10\,dB the same eight models span 55 points of accuracy, a range the clean benchmark hides completely, and the ranking is weakly \emph{anti}-correlated with clean accuracy. Finally, a one-line change to the training loop---draw a random SNR per mini-batch---recovers most of the lost accuracy for every architecture we tried, at the cost of a few points on the clean set.

This paper makes five contributions:
\begin{enumerate}
\item \textbf{A saturation audit.} Under a leakage-safe protocol (recording-level split, zero shared sampling points, verified by hashing), eight lightweight models are shown to be statistically indistinguishable on clean CWRU data, and the inflation attributable to overlapping-window leakage is bounded at $\leq 0.6$ points at full data and $\leq 0.9$ points in a 20-windows-per-class regime---an order of magnitude smaller than commonly assumed.
\item \textbf{A noise-stress matrix.} The same models, trained identically, differ by up to 55 points at 10\,dB SNR while differing by 3 points on clean data; the clean ranking does not predict the noisy ranking (Spearman $\rho=-0.38$, $n=8$).
\item \textbf{An attribution study.} By inserting a wide-kernel front-end into narrow-kernel backbones and by adding per-batch noise augmentation, we separate architectural from training-recipe effects, and show that the recipe dominates while the two are complementary.
\item \textbf{An open-set audit.} With three of ten classes withheld, closed-set accuracy stays near 100\% while rejection quality varies from 0.12 to 0.87 TPR at 5\% FPR across architectures, and it is anti-correlated with noise robustness ($\rho=-0.64$)---a deployment risk that closed-set tables cannot show. Adding the augmentation that repairs noise accuracy does not repair rejection: mean AUROC falls from 0.90 to 0.80, so the two objectives trade off and must be decided at the decision layer.
\item \textbf{A reproducible protocol and reporting checklist} (Appendix~A), including the exact split files, per-seed results and a seven-item list of information we believe every diagnostic paper should report.
\end{enumerate}
We are explicit about what this paper is \emph{not}: it proposes no new state-of-the-art architecture, and all experiments use public datasets. Its claim is narrower and, we hope, more durable: \emph{clean accuracy is not a proxy for diagnostic capability, and the community's standard evaluation cannot see the difference.}

\section{Related work}

\paragraph{Lightweight 1D architectures.} WDCNN \citep{zhang2017wdcnn} established the wide-first-kernel design that most subsequent work follows: a first layer with a kernel of 64 samples and a stride of 16 acts as a learnable band-pass filter and reduces temporal resolution before the deep stack. TICNN \citep{zhang2018ticnn} adds ensemble training and dropout; MCNN \citep{zhao2019mcnn} uses multi-scale kernels; ResNet1D, DS-CNN and LeNet1D are imported from the general 1D-CNN and keyword-spotting literature \citep{he2016resnet,hannun2019dspeech}. More recent work replaces convolution with attention or mixing blocks, e.g.\ ViT-style \citep{dosovitskiy2021vit} and MLP-Mixer-style \citep{tolstikhin2021mixer} backbones, or with graph, contrastive or domain-adversarial machinery for cross-domain diagnosis \citep{tim2025a,eaa2025a}. We include one representative of each family so that conclusions are not tied to a single design lineage.

\paragraph{Noise robustness.} Adding white Gaussian noise at test time is the field's standard robustness probe, but protocols differ widely: some papers train on clean data only and test from $-4$ to 10\,dB, others train with noise, others use coloured, harmonic or impulse noise \citep{mst2024a,mssp2025a,meas2023a}. Reported robustness therefore cannot be compared across papers, and a model robust to white noise may fail on the impulsive noise of a real gearbox. We fix one noise family for the main matrix and use other families as a robustness check.

\paragraph{Evaluation methodology.} Critiques of benchmark practice exist in adjacent fields---data leakage in time-series classification, shortcut features in image benchmarks, reproducibility in machine learning at large. For vibration diagnosis specifically, recent position papers flag overlapping-window leakage and non-comparable splits \citep{phme2026,ssrn2024}, and review articles note that reported accuracy has saturated \citep{aip2025a}. We are not aware of a study that quantifies the effect across a common architecture set under a single harness; our audit is designed to be reusable, since it needs only hashes of the window index sets.

\paragraph{Open-set diagnosis.} Most papers assume the test set contains only fault classes seen in training. In practice a machine can develop a fault that was never labelled, and a confident wrong answer is worse than an abstention. Open-set recognition is well studied in vision but rare in vibration diagnosis, and where it appears it is evaluated with different metrics. We report AUROC and TPR at a fixed 5\% false-acceptance rate.

\section{Protocol}

\subsection{Data}
CWRU \citep{smith2015cwru}: 40 recordings (10 classes: normal plus three fault diameters $\times$ three fault locations, 0.007--0.021\,in), four load conditions (0--3\,hp), 12\,kHz drive-end accelerometer, 1024-sample windows with 50\% overlap. Ten classes are used throughout.

Paderborn \citep{lessmeier2016paderborn} is used as a second dataset: six healthy and twelve artificially damaged bearing settings at 64\,kHz, 1{,}439 records of 4\,s each, aggregated into two classes (healthy, outer race) with the same windowing, split and training logic. The twelve damaged bearings in the released subset all carry outer-race damage; inner-race settings exist in the database but are not part of our subset, so the Paderborn experiment tests generalisation of the \emph{evaluation method} rather than a three-class diagnosis problem.

\subsection{Splits}
Four protocols, all released as index files:
\begin{itemize}
\item \textbf{leaky}: pool all windows, stratified random 80/20 split by class. Overlapping windows share samples across the split.
\item \textbf{safe}: for each recording the first 70\% of windows go to training and the last 30\% to test, with a half-window guard band; no sampling point appears in both sets (verified by hashing the sample-index sets).
\item \textbf{crossload}: train on 0 and 1\,hp, test on 2 and 3\,hp.
\item \textbf{filedisjoint}: 25\% of the recordings per class are held out entirely.
\end{itemize}

\subsection{Models and training}
Eight architectures, 3{,}018--747{,}722 parameters, trained from scratch with AdamW (lr $=2\times10^{-3}$, weight decay $10^{-4}$), cosine schedule, batch 128, 20 epochs, three seeds, per-window $z$-score normalisation, no pre-training. All models use the same optimiser budget, so differences reflect architecture, not tuning effort. Training runs on Apple MPS; a single run takes 8--84\,s.

\subsection{Noise and augmentation}
Test noise is additive white Gaussian at SNR $\in\{15,10,5,0,-5\}$\,dB, applied to the raw window before normalisation; “clean” corresponds to the uncorrupted window. Clean-trained models see no noise during training; augmented models draw $\mathrm{SNR}\sim U[0,15]$\,dB independently per mini-batch.

To test whether the conclusions depend on the noise family, we additionally corrupt the test set with three deterministic families at the same nominal SNR: (i) \emph{harmonic} interference, a sinusoid at a per-window frequency $f\in[0.02,0.40]$ cycles/sample plus its $2.3f$ harmonic; (ii) \emph{impulsive} noise, 16 impulses of amplitude 3 placed at uniformly random sample positions (one per 64 samples); (iii) \emph{drift}, a single low-frequency sinusoid with $f\in[0.002,0.02]$ cycles/sample and random phase. Each realisation is rescaled to the requested SNR with respect to the signal power of that window, so the four families are compared at equal energy.

\subsection{Open-set protocol}
Three classes (the 0.021-in faults at three locations) are removed from training; the model is trained on the remaining seven and evaluated on a test set containing all ten. The softmax maximum is the confidence score. We report AUROC, TPR at 5\% FPR on unknowns, and accuracy on knowns when 5\% of knowns are rejected.

\subsection{Statistics}
Means and standard deviations over three seeds. Paired comparisons use the Wilcoxon signed-rank test over the $8\times3=24$ (model, seed) pairs with Holm correction. Rank agreement between clean and noisy rankings uses Spearman's $\rho$. Every number in this paper is regenerated from \texttt{results/runs/} by one command.

